\chapter{Static parallel \hp-adaptive \glsfmtlongpl{fem}}
\label{ch:parallel}
\glsresetall

Any kind of numerical method requires thorough thought on designing suitable algorithms and data structures with respect to correctness, robustness and performance. In general, the ideas behind implementations of these methods are similar and can be generalized.
%, especially when they are based on the same numerical method.
This is also the case for additional enhancements that improve basic realizations. In case of the \gls{fem} such features are, for example, parallelization, adaptive methods, continuous or discontinuous Galerkin methods, and the support for complex geometries.

In this chapter, we will present algorithms and data structures for parallel, \hp-adaptive \glsfmtlongpl{fem}. For now, we focus on static domains with a fixed grid resolution and distribution of finite elements that will not change during the course of the application. Those can be used for multi-physics applications, in which different characteristics can be assigned to certain parts of the domain using dedicated finite elements.

Generalized thoughts on each of the two aspects have already been presented: \textcite{bangerth2009} developed a general implementation for \hp-adaptive \gls{fem} software; and a generalized distributed computing approach of the \glsfmtlong{fem} has been introduced by \textcite{bangerth2012}. However, the consolidation of both features is not trivial, thus this chapter should be understood as an enhancement of the two aforementioned contributions and bases on large parts on it. We recommend a previous reading of both articles.

We will elaborate on the trickiest facets in the following sections, after providing a brief overview about the basics that are necessary for its understanding. A development log on all features added to the \dealii{} library in the context of \hp-adaptive \gls{fem} can be found here \textcite{dealiiissue3511}.



%\setcounter{section}{-1}
\section{Prerequisities}
\label{sec:prerequisities}

% \subsection{Basics of the \glsfmtlong{fem}}

The basic idea of the \glsfmtlong{fem} conforms to the specification of a function space and finding a solution to the investigated problem in it. To be more precise, we pick a suitable set of basis functions for which the solution is a linear combination. Its coefficients are called unknowns or \glspl{dof}, since their values are determined after solving the problem.

In general, \gls{fem} requires a subdivision of the domain $\Omega$ into smaller cells $K$, where $\Omega = \cup_{i} K_i$ and each cell $K$ corresponds to a connected Lipschitz domain with nonempty interior \todo{Check Ciarlet definition?}. Typical choices for cells are triangles or quadrilaterals in two and tetrahedra or hexahedra in three dimensions.
%All these cells are mappings of a reference cell, to which we assign a set of shape functions with corresponding support points.
%Those functions are designed to have the value one on exactly one support point and the value zero on all others.
Further, on each cell $K$, we define a finite-dimensional function space $P$, accompanied by a set of functionals $N = {\Psi_1, \dots, \Psi_n}$ that are a basis to its dual space $P'$, where $n$ corresponds the number of nodes. The tuple $(K,P,N)$ is called a finite element. \todo{Cite Ciarlet or Brenner. Depends on definition of K.}

With this definition, we are able to find a set of basis functions ${\varphi_1, \dots, \varphi_n}$ of $P$ dual to $N$ fulfilling $\Psi_i(\varphi_j) = \delta_{ij}$. In other words, these basis functions span the function space $P$ while each of them has the value one on their associated node and zero on the other. They form our finite element approximation as a linear combination $u_{\text{hp}}(\vec{x}) = \sum_{j} U_j \, \varphi_j(\vec{x})$, where coefficients $U_j$ are known as \glsfmtlongpl{dof}.

We transform the original problem into a variation one by converting it to a weak formulation. This process belongs to the class of Galerkin methods, of which we distinguish between \gls{cg} and \gls{dg} methods. For the former, we require that our finite element approximation is continuous across cell interfaces resulting in shared \glspl{dof} on all transitions. For discontinuous methods, jumps of the approximation are possible on cell interfaces, and thus each cell has its own set of \glspl{dof} from which none are shared. Here, relations between cells are quantified via penalty methods that correlate \glspl{dof} on interfaces of neighboring cells.

With the weak formulation of the investigated problem, we are left to assemble and solve the corresponding system of linear equations for the unknown coefficients or \glspl{dof}. In practice, we calculate the corresponding integrals using quadrature rules and basis functions which have been mapped from a reference cell $\widehat{K}$.

When dealing with finite elements in practice, it is sufficient to work with nodes and shape functions once the type of element is specified. One of the most common types of finite elements are Lagrangian elements, that are based on Lagrangian interpolation. The distribution of nodes on quadrilaterals in Lagrangian elements up to fourth order polynomials are depicted in Fig.~\ref{fig:lagrange}.

\begin{figure}
\begin{subfigure}{.24\textwidth}
  \centering
  \begin{tikzpicture}[scale=2]
\def\Length{1}
\def\Radius{0.03}

\LagrangeCell{0}{0}{\Length}{\Radius}{1}
{{,,,}};
\end{tikzpicture}
  \caption{$Q_1$ element.}
\end{subfigure}
\begin{subfigure}{.24\textwidth}
  \centering
  \begin{tikzpicture}[scale=2]
\def\Length{1}
\def\Radius{0.03}

\LagrangeCell{0}{0}{\Length}{\Radius}{2}
{{,,,,,,,,}};
\end{tikzpicture}
  \caption{$Q_2$ element.}
\end{subfigure}
\begin{subfigure}{.24\textwidth}
  \centering
  \begin{tikzpicture}[scale=2]
\def\Length{1}
\def\Radius{0.03}

\LagrangeCell{0}{0}{\Length}{\Radius}{3}
{{,,,,,,,,,,,,,,,}};
\end{tikzpicture}
  \caption{$Q_3$ element.}
\end{subfigure}
\begin{subfigure}{.24\textwidth}
  \centering
  \begin{tikzpicture}[scale=2]
\def\Length{1}
\def\Radius{0.03}

\LagrangeCell{0}{0}{\Length}{\Radius}{4}
{{,,,,,,,,,,,,,,,,,,,,,,,,}};
\end{tikzpicture}
  \caption{$Q_4$ element.}
\end{subfigure}
\caption[Lagrangian elements $Q_p$.]{First four Lagrangian elements $Q_p$ on quadrilateral unit cells with $(1+p)^\text{dim}$ nodes each.}
\label{fig:lagrange}
\end{figure}

This is just a brief introduction to \gls{fem} in order to provide the fundamentals for this chapter and to specify the nomenclature used. More details on this topic can be found in common literature \parencite[e.g.][]{girault1986, elman2014, kuzmin2015} \todo{add more literature?}. Especially \textcite{brenner2008} provided a more rigorous and mathematically sound definition on finite elements.


\subsection{Adaptive methods for \glsfmtshort{fem}}

In computational applications, adaptive methods are used to align the resolution or rather the processing power to the complexity of the investigated problem, i.e.\@ to specific sections of the domain. In terms of grid refinement, cells will be split twice per dimension, resulting in two, four or eight children in one, two or three dimensions, respectively. On the contrary, we join the corresponding amount of cells to their parent cell for grid coarsening. This process is also known as \h-adaptation, referring to adjusting the cell's length or diameter \(h\) locally. We require hierarchic relations between parent cells and their children that motivate the use of tree structures rooting in each cell of the initial coarse mesh. Combined this leads to a forest structure, in which each level corresponds to a full representation of the domain. The entirety of all leaves in this forest resembles the computational domain relevant for assembling and solving the system of equations. These particular cells, which have no children, are called active cells and are the only ones carrying \glspl{dof}. For more information, \textcite{bangerth2003} have elaborated more rigorously on grid adaptation, especially for the \gls{fem} method.

As an alternative to modifying the grid resolution, we can also adapt the function space using various finite elements associated with each cell. These finite elements may differ in the polynomial degree $p$ of their shape functions, offering the unique possibility for \p-adaptation, or \hp-adaptation when used together with grid adaptation. In practice, we specify a collection of finite elements, one of which is assigned to each cell in the domain. We identify the currently assigned finite element with an \textit{active finite element index} on each cell.

Refinement level differences on quadrilateral or hexahedral meshes, as well as varying finite elements on neighboring active cells lead to hanging nodes which are nodes with no counterpart on the opposite side of the interface. A depiction of \h- and \p-adaptive domains along with hanging nodes is presented in Fig.~\ref{fig:hpadaptivity}.

\begin{figure}
\begin{subfigure}{0.49\textwidth}
  \centering
  \begin{tikzpicture}[scale=2]
\def\Length{0.5}
\def\Radius{0.03}

\LagrangeCell{-2*\Length}{0}{2*\Length}{\Radius}{2}
{{,,,,,,,,}};

% remove face node which becomes a hanging node on a vertex
\draw[fill=white,draw=white] (0,\Length-\Radius) rectangle (-\Radius,\Length+\Radius);

\LagrangeCell{0}{0}{\Length}{\Radius}{2}
{{,,,,,,,,}};
\LagrangeCell{\Length}{0}{\Length}{\Radius}{2}
{{,,,,,,,,}};
\LagrangeCell{0}{\Length}{\Length}{\Radius}{2}
{{,,,,,,,,}};
\LagrangeCell{\Length}{\Length}{\Length}{\Radius}{2}
{{,,,,,,,,}};
\end{tikzpicture}
  \caption{\h-adaptation.}
  \label{fig:hapaptivity}
\end{subfigure}
\begin{subfigure}{0.49\textwidth}
  \centering
  \begin{tikzpicture}[scale=2]
\def\Length{1}
\def\Radius{0.03}

\LagrangeCell{0}{0}{\Length}{\Radius}{2}
{{,,,,,,,,}};
\LagrangeCell{\Length}{0}{\Length}{\Radius}{4}
{{,,,,,,,,,,,,,,,,,,,,,,,,}};
\end{tikzpicture}
  \caption{\p-adaptation.}
  \label{fig:papaptivity}
\end{subfigure}
\caption{Example of an \h-adapted mesh consisting of quadrilateral cells. Differences in the refinement level of neighboring cells gives rise to hanging nodes. Note that both scenarios have the same number of \glspl{dof}.}
\label{fig:hpadaptivity}
\end{figure}

In combination with \gls{cg} methods, the finite element approximation needs to be continuous on these hanging nodes, which requires constraining them to the surrounding ones on the shared interface.

For \h-adaptive methods, we need to constrain hanging nodes from cells of a finer refinement level to the neighboring coarser cell by interpolation. A typical situation is depicted in Fig.~\ref{fig:hapaptivity}. For convenience, the level difference of neighboring cells is usually limited to one level in order to simplify interpolation between cell interfaces. For quadrilateral or hexahedral meshes, this \textit{2:1 mesh balance} ensures to limit the occurrence of hanging nodes.

However, when different finite elements face each other in case of \p-adaptation, we need to impose continuity on these interfaces as well. This is performed by restricting the element with the higher polynomial degree to the continuity of the lower one. Following the example from Fig.~\ref{fig:papaptivity}, we would have to restrict the additional nodes of the $Q_4$ element to the ones of the $Q_2$ element on their shared line. To formulate it in a more general way, we need to constrain continuity of all neighboring finite elements to their largest common finite element subspace. We say that this particular element \textit{dominates} the others.

There may be cases in which neighboring elements do not dominate each other, since they do not pose compatible continuity requirements. In this case, we pick a common subspace from the full collection. As an example, this would be the case for two neighboring vector-valued elements $(Q_2 \times Q_1)$ and $(Q_1 \times Q_2)$, which would be dominated by a $(Q_1 \times Q_1)$ element.

This summary recaps data structures and algorithms for \hp-adaptive \gls{fem} which have been elaborated in great detail by \textcite{bangerth2009}.

%This summary recaps data structures and algorithms for \hp-adaptive \gls{fem} as well as parallel \h-adaptive \gls{fem} which have been elaborated in great detail by \textcites{bangerth2012}{bangerth2009}, respectively.

% make this part of the next chapter?
%When the function space is a subset of a different finite element, we say that it dominates the other finite element.

%Dominating finite elements own \glspl{dof} on cell interfaces, with some minor exceptions as elaborated by \textcite{bangerth2009} who also provide more information on \hp-refinement.


\subsection{Parallelization of \glsfmtshort{fem}}

\Gls{hpc} applications requires scalable algorithms and data structures for machines with distributed memory, which we realize using the \gls{mpi} standard. In this context, the workload on all participating processes will be shared by partitioning the whole domain among them.

These ambitions pose different challenges on the design of data structures and information exchange between processes. Especially dynamic changes of the domain by \hp-adaptation are difficult, since they require the redistribution of the workload by repartitioning the global mesh, as well as the transfer of data not only between processes, but also between the former and the updated mesh. Here, we distinguish between different subsets of cells and \glspl{dof} which will be introduced in the following.

Each process will only store data of cells that belong to its owned section of the domain, which we call \textit{subdomain}. Those cells are called \textit{locally owned cells}, and all \glspl{dof} on them are referred to as \textit{locally active \glspl{dof}}. Every \gls{dof} will be uniquely enumerated on the global mesh. With \gls{cg} methods, \glspl{dof} on interfaces between cells of different processes may be owned by either one or the other. Thus, not all \glspl{dof} on locally owned cells have to be \textit{locally owned \glspl{dof}}.

During the assembly of equation systems, we need to refer to surrounding cells that do not belong to the local domain. Typically, the halo spanning one level of cells around the locally owned domain is sufficient. We call these cells \textit{ghost cells}. Data from ghost cells has to be requested from the owning process. The combination of locally owned and ghost cells and their corresponding \glspl{dof} are labeled \textit{locally relevant}.

Further, we require that every process stores a copy of the common coarse mesh. This allows an easy construction of the whole grid during repartitioning just with adaptive methods. We establish the 2:1 mesh balance on all cells of the local mesh, leaving refined cells that are not actually used. These cells that are not locally relevant will be called \textit{artificial cells}. An example of a distributed, adapted mesh is shown in Fig.~\ref{fig:paralleldistribution} with all presented classes of cells.

\begin{figure}
\begin{subfigure}[t]{.35\textwidth}
  \centering
  \begin{tikzpicture}
\def\Length{1}
\def\Radius{0.03}

\LagrangeCell{0}{0}{\Length}{\Radius}{4}
{{,,,,,,,,,,,,,,,,,,,,,,,,}};
\end{tikzpicture}
  \caption{Process with rank 0.}
\end{subfigure}
\begin{subfigure}[t]{.35\textwidth}
  \centering
  \begin{tikzpicture}
\def\Length{1}
\def\Radius{0.03}

\LagrangeCell{0}{0}{\Length}{\Radius}{4}
{{,,,,,,,,,,,,,,,,,,,,,,,,}};
\end{tikzpicture}
  \caption{Process with rank 1.}
\end{subfigure}
\begin{subfigure}[t]{.28\textwidth}
  \centering
  \begin{tikzpicture}
\begin{customlegend}[legend entries={locally owned,ghost,artificial,locally relevant}]
\addlegendimage{fill=red,area legend};
\addlegendimage{fill=yellow,area legend};
\addlegendimage{fill=green,area legend};
\addlegendimage{black,dashed,sharp plot};
\end{customlegend}
\end{tikzpicture}
\end{subfigure}
\caption{Parallel distribution of an \h-adapted mesh on two processes.}
\label{fig:paralleldistribution}
\end{figure}

%This is just a brief outline of all the requirements that \textcite{bangerth2012} worked out, which is crucial for the upcoming section.

%We work with \dealii{}. This library works with quadrilaterals and hexahedrals only. All the considerations in this dissertation should be easily applicable on triangles and tetrahedra as well.

%In the \dealii{}, the parallel is handed out to an oracle that is requested on all mesh related operations. \dealii{} offers an interface to the \pforest{} library.

\textcite{bangerth2012} worked out the parallelization of \h-adaptive \gls{fem} within the \dealii{} library in great detail.
\section{Enumeration of \glsfmtlongpl{dof}}
\label{sec:enumeration}

Formulating and solving the system of linear equations requires an unique identification of all involved \glspl{dof} in the global mesh.

\Glspl{dof} are associated with mesh objects, i.e., vertices, edges, faces, and cells. If support points are located on interfaces between neighboring cells in the context of \gls{dg} methods, they are assigned separate \glspl{dof} on each cell. Thus, the enumeration of both \glspl{dof} and cells can happen analogously. However, \gls{cg} methods require that shared \glspl{dof} on interfaces between neighboring cells are unique. Thus, each \gls{dof} has to relate to a single cell, or in other words, will be owned by a single cell. This assignment is crucial for the preparation of relevant distributed data structures.

%Depending on which Galerkin method has been picked, we need to ... .
%Continous Galerkin methods require a unique identification of \glspl{dof} on interfaces between cells, because of the continuity requirement on the solution.

The latter scenario poses challenges in the enumeration of \glspl{dof} when considering either parallelization or \hp-adaptive methods, let alone a combination of both. A first attempt to cope with this problem would involve so called constraints: We enumerate all \glspl{dof} on each cell independently, but identify two similar \glspl{dof} as identical during the assembly of the equation systems. Although this approach would be easy to implement, we would needlessly add \gls{dof} duplicates to the equation system, sacrificing performance by wasting memory and computing time. We conclude that a unique enumeration of \glspl{dof} is mandatory for a robust \gls{fem} implementation for \gls{cg} methods.

Both \textcite{bangerth2012} and \textcite{bangerth2009} have dealt with \gls{dof} enumeration with parallelization and \hp-adaptive methods, respectively, and presented algorithms for each case, but the combination of both is not trivial. In this section, we will briefly summarize each implementation and present an enhanced algorithm in detail for the unique identification of \glspl{dof} for \gls{cg} methods with parallel \hp-adaptive methods.

In the following, all algorithms will be presented independently of the underlying data structures and should be easily applicable on any kind of existing \gls{fem} code. Results will be presented afterwards with our reference implementation in the \dealii{} library.

%A common approach to parallelization for cell based methods on distributed memory machines is to assign multiple cells to one process of that form a subdomain. Typically during the assembly of the equation system we need to the values of the surrounding cells. Thus, an extension of the subdomain by so called ghosts cells is required.

%When using continuous Galerkin methods on the other hand, \glspl{dof} are shares along interfaces between cells.

%However, a different way of enumerating \glspl{dof} has to be taken with parallel \gls{fem}. \textcite{bangerth2012} provided a

%For \hp-adaptive methods, cells with different finite elements assigned may neighbor each other. Thus, we may also encounter hanging nodes on neighboring cells as depicted in Fig.~\ref{fig:papaptivity} in addition to the ones arising from \h-adaptation.

%In the context of \gls{cg} methods, these hanging nodes from \p-adaptation need to comply to the continuity condition along the faces of a cell.

%For \hp-adaptive methods, cells with different finite elements assigned may neighbor each other. Thus, we may also encounter hanging nodes on neighboring cells as depicted in Fig.~\ref{fig:papaptivity} in addition to the ones arising from \h-adaptation.

%We will also encounter shared dofs on certain neighboring finite elements.

%\textcite{bangerth2009} introduced efficient data structures for \hp-adaptive \gls{fem} codes. We distinguish between objects on the domain to which nodes and \glspl{dof} are assigned, i.e.\@ vertices, lines and quads, and hexes, depending on the dimension of the problem. For each of these object on the domain, a linked list stores the indices of all attached cells and their corresponding \glspl{dof}.

%Problems that occur involve shared \glspl{dof} on neighboring cells with the same finite element assigned on either the same or different subdomains, as well as with different finite elements assigned on either the same or different subdomains.

For \hp-adaptive \gls{fem}, the algorithm proposed by \textcite[Sec.~4.2]{bangerth2009} enumerates all \glspl{dof} on each cell consecutively in a first step, and then unifies these shared \glspl{dof} on cell interfaces by favoring the index from the dominating finite element.

%Enumerate degrees of freedom on each cell consecutively. Different indices for degrees of freedom will be assigned on interfaces with neighboring cells. A linked list is introduced for each lower–dimensional object (i.e. vertices, lines, faces), which stores the indices of all degrees of freedom located on it.

%Unify shared degrees of freedom on interfaces between cells. Depending on the finite elements used on neighboring cells, we identify overlapping degrees of freedom with the previously introduced lists, and merge them by keeping the lower index and eliminating the other.

%Later, they indices of dof duplicates will be unified using finite element, if both finite elements are compatible?

%Since workload is assigned to cells, we will assign \glspl{dof} to cells with a lower number of dofs. To distinguish between cells, we employ finite element domination algorithm.

In parallel applications, the enumeration of \glspl{dof} on interfaces between neighboring subdomains pose a problem: They have to be assigned to one of them, for which \textcite[Sec.~3.1]{bangerth2012} proposed to use a certain \textit{tie-break} criterion as a decision aid. Their algorithm starts by enumerating \glspl{dof} on all locally owned cells. On interfaces between subdomains, all \glspl{dof} will be assigned to the process with the lower rank and thus keep the index from the subdomain with the lower identifier. Once ownership of all \glspl{dof} is clarified, they will be re-enumerated on the global domain so that every locally owned \gls{dof} has its final index assigned. Each process needs to know all locally relevant \glspl{dof} for the solution of the equation system, which requires exchanging \gls{dof} indices on ghost cells via point-to-point communication. This operation has to be performed twice since there may have been \glspl{dof} on ghost cells of which the owning process did not know the correct indices of after the first exchange.

%The enumeration of dofs faces fundamental problems in both parallelisation and \hp-adaptive methods individually. Let us first elaborate on those and quickly recall the corresponding papers, before moving to the actual consolidation of the two.

For parallel \hp-adaptive \gls{fem}, the mere concatenation of both algorithms does not suffice. The case in which distinct finite elements from different subdomains are adjacent has to be given special consideration. We could cope with this situation by constraining these particular \glspl{dof} to each other. However, this would leave unnecessary \gls{dof} duplicates in the equation system. Additionally, the global number of \glspl{dof} would differ with the number of subdomains in this approach. We would rather keep it independent from the number of processes, which would simplify debugging and assures that our solvers yield the same results on any amount of subdomains. The algorithm developed in this thesis meets this requirement and combines the ideas of both previous algorithms.

To follow the same nomenclature as \cite{bangerth2012}, we call the set of all locally owned cells $\mathbb{T}^p_{loc}$, the set of all ghost cells $\cellsp{ghost}$, and the set of all locally relevant cells $\cellsp{rel} = \cellsp{loc} \cup \cellsp{ghost}$ on processes and subdomains identified by the integer $p$.

We need to enumerate all \glspl{dof} on locally relevant cells, which includes ghost cells. Thus, we begin by exchanging active finite element indices on ghost cells so that we know about all locally relevant finite elements and can prepare all data structures accordingly. We do this with point-to-point communication.

In short, our algorithm is based on the parallel algorithm by \textcite{bangerth2012} for the most part. We will add a \gls{dof} unification step after enumerating all locally owned cells, while subjecting to the finite element domination hierarchy to decide ownership on all interfaces. After exchanging \glspl{dof} on all ghost cells, we are left to merge the valid \glspl{dof} on interfaces with the valid counterparts.

In detail, the complete algorithm consists of the following six steps, starting with an initialization step flagged with `0`.
\begin{enumerate}
  \item[0.] \textit{Initialization.}
  On all locally relevant cells $K \in \cellsp{rel}$, \gls{dof} indices are set to an invalid value, for example $i \coloneqq -1$.
  \item \textit{Local enumeration.}
  Iterate over all locally owned cells $K \in \cellsp{loc}$ and assign valid \gls{dof} indices in ascending order, starting from zero.
  \item \textit{Tie-break.}
  Go over all locally owned cells $K \in \cellsp{loc}$. If $K$ neighbors a ghost cell which has the same finite element assigned and belongs to a subdomain of lower rank $q < p$, then invalidate all \glspl{dof} on their shared interface by setting their index to the invalid value $i$.
  \item \textit{Unification.}
  Go over all locally owned cells $K \in \cellsp{loc}$. For all shared \glspl{dof} on interfaces to neighboring cells, constrain the \gls{dof} of the dominated finite element to the one of the dominating element. If no dominating \gls{dof} can be determined, constrain the \gls{dof} with the higher index to the lower one. On interfaces to ghost cells, constrain \glspl{dof} to the invalid value $i$ if the dominating element is assigned to the ghost cell.
%  Note that at this stage, we will never overwrite a valid index on ghost interfaces.
  Populate a list with pairs of identifiers for these constrained \glspl{dof}.
%  This has to happen after the tie-break.
\end{enumerate}
At this stage, each process knows which \glspl{dof} are owned by them.
\begin{enumerate}[resume]
  \item \textit{Global re-enumeration.}
  Iterate over all locally owned cells $K \in \cellsp{loc}$
  and enumerate those \gls{dof} indices in ascending order that have a valid value assigned, while considering all constraints from the previous step. Store the number of all valid \gls{dof} indices on this subdomain as $n_p$. In a next step, shift all indices by the number of \glspl{dof} that are owned by all processors of lower rank $q < p$, or in other words, by $\sum_{q=0}^{p-1} n_q$. This corresponds to a prefix sum or exclusive scan, and can be obtained via \texttt{MPI\_Exscan} \textcite{mpi31}.
\end{enumerate}
At this stage, all subdomains and processes have the correct indices assigned to all locally owned \glspl{dof}, which are enumerated consecutively.
\begin{enumerate}[resume]
  \item \textit{Ghost exchange.}
  Communicate \gls{dof} indices from locally owned cells $K \in \cellsp{loc}$ to other processes using point-to-point communication as follows.
  \begin{enumerate}[label=\alph*.]
%    \item Mark all vertices on locally owned cells $\cellsp{loc}$.
%    \item Go over all vertices on ghost cells $\cellsp{ghost}$ and store all marked vertices and their all processor identifiers.
%    \item Mark all locally owned cells that have ghost neighbors.
    \item Prepare containers with data to be sent from subdomain $p$ to each adjacent subdomain $q$.
    \item Loop over all locally owned cells that have ghost neighbors. Add the cell identifier with all associated \gls{dof} indices to every data container that corresponds to an adjacent subdomain of the current cell.
    \item Send each data container to its designated destination process with non-blocking point-to-point communication via \texttt{MPI\_Isend} \textcite{mpi31}.
    \item Receive data containers from processes of adjacent subdomains via non-blocking point-to-point communication via \texttt{MPI\_Irecv} \textcite{mpi31}. The received data corresponds to ghost cells of this subdomain $p$. On each of these cells, set the received \gls{dof} indices accordingly.
  \end{enumerate}
  All communication in this step is symmetric, which means that a process only receives data from another process when it also sends data to it. Thus, there is no need to negotiate communication.
  \item \textit{Merge on interfaces.}
  After the previous ghost exchange each \gls{dof} on interfaces with ghost cells has exactly one valid index assigned. Go over all locally relevant cells $K \in \cellsp{rel}$. On interfaces between locally owned and ghost cells, set all remaining invalid \gls{dof} indices with the corresponding valid one of the dominating finite element.
\end{enumerate}
At this stage, all locally owned cells $K \in \cellsp{loc}$ have their correct \gls{dof} indices set.
\begin{enumerate}[resume]
  \item \textit{Ghost exchange.}
  Some ghost cells may still have invalid \gls{dof} indices assigned since the owning process did not know all correct indices on this particular cell when the last communication happened. Another ghost exchange closes this gap by repeating step 5. However this time, only data from those cells has to be communicated which had invalid \gls{dof} indices prior to step 5d.
\end{enumerate}
One variant of this algorithm would be to modify steps 2 and 3 inasmuch as we apply the tie-break criterion on all \glspl{dof} on ghost interfaces and exclude them from \gls{dof} unification entirely. However, this would not assign shared \glspl{dof} to the dominating finite element on ghost interfaces, which would be inconsistent to the locally owned parts of the domain.

At the end of this algorithm, all global \gls{dof} indices have been set correctly, and every process knows the indices of all locally relevant \glspl{dof}: All \glspl{dof} on interfaces belong to the dominating finite element on the process with the smallest rank. No new communication steps in addition to the two original ones from the \h-adaptive variant are introduced.

A suitable benchmark that we used to test the enumeration algorithm consists of a two-dimensional domain of four neighboring cells. We provide two different Lagrangian elements that share at least one additional \gls{dof} per cell interface than only on vertices. For this purpose, we choose Lagrangian elements $Q_2$ and $Q_4$. Each of these finite elements will be assigned to two catty-cornered cells. Furthermore, we divide the mesh into two subdomains containing two neighboring cells of differing finite elements. The whole setup is shown in Fig.~\ref{fig:enumbenchmark} and covers all combinations of adjacent finite elements that we have encountered so far in the parallel \hp-adaptive context. A step-by-step demonstration of the algorithm on this particular benchmark is presented in App.~\ref{app::enumeration}.

\begin{figure}
\centering
\begin{tikzpicture}[scale=3]
\def\Length{1}
\def\Radius{0.03}

\fill[color=green] (0,0) rectangle (2*\Length,\Length);
\fill[color=yellow] (0,\Length) rectangle (2*\Length,2*\Length);

\LagrangeCell{0}{0}{\Length}{\Radius}{2}
  {{0,1,2,3,4,5,6,7,8}};
\LagrangeCell{\Length}{0}{\Length}{\Radius}{4}
  {{1,9,3,49,10,5,11,12,13,14,15,16,17,18,54,19,20,21,22,23,24,25,26,27,28}};
\LagrangeCell{0}{\Length}{\Length}{\Radius}{4}
  {{2,3,29,50,30,31,32,33,52,34,35,7,36,37,38,39,40,41,42,43,44,45,46,47,48}};
\LagrangeCell{\Length}{\Length}{\Length}{\Radius}{2}
  {{3,49,50,51,52,53,54,55,56}};
\end{tikzpicture}
\begin{tikzpicture}
\begin{customlegend}[legend cell align=left]
\addlegendimage{fill=green,area legend};
\addlegendentry{subdomain 0};

\addlegendimage{fill=yellow,area legend};
\addlegendentry{subdomain 1};
\end{customlegend}
\end{tikzpicture}
\caption[Benchmark scenario for \gls{dof} enumeration.]{Benchmark scenario to verify our algorithm for \gls{dof} enumeration. The mesh is distributed on two \gls{mpi} processes, each owning one $Q_2$ and one $Q_4$ element. All \glspl{dof} are uniquely identified on the global mesh as a result of the enumeration algorithm from Sec.~\ref{sec:enumeration}.}
\label{fig:enumbenchmark}
\end{figure}

There might be chains of constraints that span over \glspl{dof} from multiple subdomains. To deal with this case, we might need more than the current two communication steps in our algorithm. However, we could not think of any scenario in which this is going to be the case, and did not encounter any issues in our investigations so far.

In three-dimensional scenarios, \textcite[Sec.~4.6]{bangerth2009} pointed out possible complications with circular constraints during \gls{dof} unification whenever three different finite elements share a common edge. We still have not figured out a satisfactory solution for this problem and subject to their original way to cope with this situation: All \glspl{dof} on these edges will be excluded from the unification step and will be treated separately via constraints.
\section{Load balancing}
\label{sec:balancing}

The efficient use of all computational resources requires a uniform distribution of all workload among them. There are many factors that determine the workload in a \gls{fem} application, above all the preparation of data structures, the assembly of both the matrix and right hand side of the linear equation system, and the choice of the type of solver. Each of these categories contributes a different amount to the total workload, with the solver accounting for the largest share in general. We should thus balance the number of rows or non-zero entries per process in the solution matrix. However, information about the matrix is available late in a solution cycle, so we need to look for a different measure at an earlier stage.

In most \h-adaptive applications, cells are similar and correspond to the same workload. Thus, we can simply balance the number of cells on all processes. However with \hp-adaptive application, this is no longer the case due to the diversity of finite elements. %Here, the workload varies with the number of \glspl{dof} on each cell. (and the quadrature)
In this case, since the domain is partitioned on the basis of cells, we need to assign a corresponding weight to every cell that determines its individual workload and balance the cumulated weights among all processes.

The workload of each cell depends on its assigned finite element and quadrature formula, and correlates to the number of \glspl{dof} and quadrature points, among other quantifiable values that depend on the individual problem. For example, Lagrangian elements of different order as depicted in Fig.~\ref{fig:lagrange} each have a distinct number of \glspl{dof}.

%For example, consider Lagrangian elements of different polynomial degrees with different number of \glspl{dof} depicted in Fig.~\ref{fig:lagrange} require a different workload, mainly consumed by matrix assembly and the solution process.

%and can be quantified
%We expect a correlation of every cell's individual workload with the assigned number of \glspl{dof} and number of quadrature point of the assigned formula, among others quantifiable values.

%Ideally, we want to balance the workload. We will thus perform weighting.

%and balance the cumulated sum of cells to be equal on all processes. We can do that with a prefix sum.

For the purpose of load balancing, \textcite[Sec.~3.3]{burstedde2011} provided an algorithm for weighted partitioning and enhanced \pforest{} \textcite{p4est22} with a corresponding implementation, from which we take advantage in \dealii{}. Omitting details about the communication between processes, we will briefly outline its basic idea: On a distributed mesh, calculate the prefix sum of cell weights in the global scope, determine the partition boundaries with a binary search, and transfer cells to their new owning processes if necessary.

%The partitioning algorithm of \pforest{}

%\pforest{} offers such a weighting mechanism. The basic idea behind it is, that we calculate the cumulated sum of cells with a prefix sum, and assigning the correpsonding partition to our process with a binary search.

%We rely on the implementation of \textcite{burstedde2018} in \pforest{} \textcite{p4est22}. The basic idea of this algorithm is to form a cumulated of all cell weights with an \texttt{MPI\_Exscan} call, and then each process will find its balanced range of cells with binary searches.

In the context of \hp-adaptive \gls{fem} applications, we identify the assembly of the linear equation system and its solution as the most expensive tasks, and correlate their individual contribution to the workload on each cell with its number of \glspl{dof} provided by the associated finite element.
%As a first indicator for the workload, we make the number of \glspl{dof} responsible.
The total workload of iterative solvers combined with multigrid preconditioning ideally scales with the global number of \glspl{dof}, i.e., $\mathcal{O}\left(N_\text{dofs}\right)$. We therefore expect that the corresponding workload of each cell attributed to the solution process also scales with the number of cell-related \glspl{dof}, i.e., $\mathcal{O}\left(n_\text{dofs}\right)$. Furthermore, we suppose that the workload of each cell for the matrix assembly will be of order $\mathcal{O}\left(n_\text{dofs}^2\right)$ since all \glspl{dof} in a cell couple with each other.

We expect that the actual workload of an \hp-adaptive \gls{fem} application will actually scale with an order somewhere in between the two, i.e.\@ $\mathcal{O}\left(n_\text{dofs}^c\right)$ with a constant exponent $c \in [1,2]$. We use this strategy for investigations in this dissertation in which we also determine a suitable exponent $c$.

Independent of the approach just described, weighting each cell with a linear combination $(a \, n_\text{dofs}^2 + b \, n_\text{dofs})$ also appears conceivable, for which the partitioning results depend on the ratio of both constants $a$ and $b$, rather than an exponent $c$.

%We will assign a weight $n_\text{dofs}^i$ to every cell, % and balance the cumulated sum of cells to be equal on all processes. We can do that with a prefix sum.

A reliable measure of weights is tied to the type of problem and the choice of the solver. With the presented approach, we still have to specify a suitable weight manually. It is part of future work to supply a heuristic from which we will determine suitable weights automatically. %, depending on the type of problem and the choice of the solver.